% ==================================================================
%  ABSTRACT + INTRODUCTION (pp. 1--2)
% ==================================================================

\begin{abstract}
Selecting among compound LLM agent designs today reduces to benchmark
racing: run every candidate end-to-end and pick the winner.
This approach conflates cost, model competence, routing quality, and
task structure into a single opaque score.
We prove an exact four-term identity that decomposes the performance gap
between any two designs into these four factors, yielding a closed-form
crossover depth $D_{\mathrm{cross}} = \log(1-p_l)/\log(1-p_s)$ at which a
cheap model with retries overtakes an expensive one.
Building on this decomposition, we give a design-selection algorithm
whose evaluation cost is independent of the number of candidate designs:
at $k{=}500$ candidates it uses $33\times$ fewer evaluations than racing.
For reusing designs across task domains, we derive a three-case transfer
rule---plug-and-play, re-route, or start over---with explicit thresholds
determined by two formally independent notions of task equivalence.
The transfer rule is conservative: its worst-case decision accuracy is
73.8\%.
\end{abstract}

\section{Introduction}
\label{sec:intro}

A practitioner building an LLM agent system faces a combinatorial design
space: which model, how many retries, what routing policy, what budget
allocation across task modes.  The standard response is benchmark
racing---evaluate every candidate end-to-end on a held-out set and pick the
winner.  Racing is model-agnostic and assumption-free, but it scales
linearly in the number of candidates and tells the practitioner nothing
about \emph{why} one design beats another.  When the winning design is
deployed on a new task domain, the practitioner must race again from
scratch.

This paper treats agent design selection as a structured optimization
problem and provides three tools---one for measurement, one for decision,
and one for reuse---that together replace opaque racing with a principled
pipeline.

\paragraph{Pillar 1: Decompose (measurement).}
Under a packed latent-mode family assumption, we prove that the
log-performance gap between any two agent designs decomposes exactly into
four terms: per-step \emph{cost}, model \emph{competence}, routing
\emph{information gain}, and routing \emph{mismatch}
(Theorem~\ref{thm:four-term}).  Each term is estimable from a small pilot
sample.  The decomposition yields a crossover formula:
$D_{\mathrm{cross}} = \log(1-p_l)/\log(1-p_s)$, where $p_l$ and $p_s$ are
the per-attempt success probabilities of a large and a small model.  When
$D_{\mathrm{cross}}$ falls below the cost ratio, the small model with
retries dominates.  The four terms are components of the primal objective,
not Lagrangian shadow prices; a natural dual interpretation is precisely
wrong (Appendix~\ref{app:lagrangian}).

\paragraph{Pillar 2: Decide (algorithm).}
The separable structure of the four-term identity means that once per-mode
success probabilities are estimated from a shared pilot, the objective for
\emph{any} design in the menu can be computed analytically---at zero
additional evaluation cost.
Algorithm~\ref{alg:decomp-guided} exploits this to select a near-optimal
design with regret at most $2\delta$ (Proposition~\ref{prop:near-opt}),
where $\delta$ is the pilot estimation accuracy.
Its evaluation cost is $O(m \cdot \Mn)$ model--mode pairs, independent of
the design-space size~$k$.
At $k{=}500$, this yields a $33\times$ reduction in evaluations relative to
racing, with comparable regret.

\paragraph{Pillar 3: Reuse (transfer).}
When a practitioner moves to a new task domain, two questions arise: can the
model selection be reused, and can the routing be reused?  We formalize
these as \emph{post-training equivalence} (same competence profiles) and
\emph{orchestration equivalence} (same routing structure), and prove that
the two notions are logically independent
(Proposition~\ref{prop:independence}).
This independence generates a three-case decision rule with quantitative
thresholds (Theorem~\ref{thm:transfer}): if both divergences are small,
transfer the design directly (plug-and-play); if only the competence
profiles match, keep the model but re-estimate routing (re-route);
otherwise, start over.  In the re-route case, the pilot cost drops by a
factor of~$m$ relative to full re-optimization.

\paragraph{Contributions.}
\begin{enumerate}
\item An exact four-term performance decomposition with a closed-form
  crossover formula, providing interpretable diagnostics for any pair of
  agent designs (Section~\ref{sec:decomposition}).
\item A decomposition-guided selection algorithm with near-optimal regret
  and evaluation cost independent of the design-space size
  (Section~\ref{sec:algorithm}).
\item A transfer theory with two independent equivalence notions and a
  three-case decision rule with explicit thresholds
  (Section~\ref{sec:transfer}).
\end{enumerate}
